% !TeX root = ../main.tex
%------------------------------------------------------------------------------
%  Ground Behaviour and Face Stability
%  Sources: Brady & Brown (2006) Ch. 7-9; Hoek & Brown (1980) Ch. 7;
%  Broms & Bennermark (1967); Davis et al. (1980); Anagnostou & Kovari (1996).
%------------------------------------------------------------------------------
\chapter{Ground Behaviour and Face Stability}
\label{ch:ground-behaviour}

Rock mass classification describes the quality of the ground but not the
manner in which it will fail. Rock masses
of equal rating may fail by quite different mechanisms, one by a block
released from the crown and another by overstressing and spalling of the same
crown, and the analysis, support and construction sequence appropriate to each
have little in common. This chapter therefore organises ground by behaviour
rather than by quality. Classification by behaviour is the older of the two ideas in
underground engineering and remains the more useful: massive, stratified and
blocky rock behave differently because different elements carry the load, and
weak ground behaves differently again because the face becomes the governing
problem. Which of these applies determines every subsequent step of the
design, and it is established from geology and stress rather than from a
rating.

\section{Ground type as the design frame}
\label{sec:ground-type-frame}

\subsection{Behaviour classes}
\label{sec:behaviour-classes}

Six behaviour classes account for most ground encountered underground, and
each is diagnosed
from a combination of intact strength, discontinuity condition and stress
level rather than from any single index.

\emph{Massive} rock is sparsely jointed and strong, so the rock mass behaves
as a continuum and the boundary stress distribution governs.
\emph{Stratified} rock is bedded or laminated, and the immediate roof behaves
as a beam or plate spanning between abutments. \emph{Blocky} rock is defined by
several persistent discontinuity sets, and behaviour is governed by the
geometry of the blocks they form. \emph{Heavily fractured} rock has
discontinuity spacing small relative to the opening, so individual blocks lose
meaning and the mass again approximates a continuum, now a weak one.
\emph{Weak and soft ground} has low intact strength, and failure occurs
through the material rather than around it, with face stability becoming the
governing problem. \emph{High-stress} ground is strong rock at stress
sufficient to induce failure at the boundary, producing spalling and, at
extremes, bursting.

The diagnostic quantities are the ratio of rock mass strength to in-situ stress
and the ratio of block size to opening span, $l_b/B$
(Figure~\ref{fig:behaviour-class-map}). Here $l_b$ is the typical edge length
of the blocks that the discontinuities form, which is set by their spacing,
and $B$ is the span of the opening. The first
distinguishes stress-driven from structurally controlled behaviour; the second
distinguishes discrete from continuum treatment. With a high strength ratio,
sparse jointing gives massive rock; several persistent sets forming blocks
that are not small relative to the opening give blocky rock; spacing small
relative to the opening gives heavily fractured rock; and where bedding or
lamination is the dominant structure the rock is stratified. A low strength
ratio gives weak and soft ground where the intact material is weak, and
high-stress ground where it is strong and brittle.

\begin{figure}[htbp]
\centering
\datafigure[\textwidth]{ch13/behaviour-class-map}
\caption{Behaviour classes on the two diagnostic ratios, with the analysis route each implies (schematic). The block-size limits of 0.1 and 1/3 are those of Section~\ref{sec:sidewalls-fractured} and apply to the upper row; the other boundaries are indicative. At a low strength ratio the class is set by the intact material, weak or strong and brittle, and not by block size. Stratified rock is identified from the type of structure, not from these ratios.}
\label{fig:behaviour-class-map}
\end{figure}

Classes change along an alignment, and difficulties concentrate at the
transitions. A drive that passes from blocky rock into a heavily fractured
fault zone changes behaviour class within a few metres, and the support
designed for one is inappropriate for the other.

\subsection{Massive rock: boundary stress and failure zone}
\label{sec:massive-rock}

In massive rock the elastic solutions of Chapter~\ref{ch:rock-mechanics}
apply, and boundary stress governs. Tangential stress concentrates at the
excavation boundary by a factor that depends on shape and on the
horizontal-to-vertical stress ratio $k$. For a circular opening in a
hydrostatic field the tangential stress at the wall is twice the field stress;
as $k$ departs from unity, stress concentrates where the boundary runs
parallel to the major principal stress and relaxes where it runs
perpendicular to it. For a circular opening the wall values are
$(3-k)\sv$ at the sidewall and $(3k-1)\sv$ at the crown
(Section~\ref{sec:kirsch-appendix}).

Where $k$ is low, the crown and invert of a circular opening can go into
tension. Since the crown value is $(3k-1)\sv$, this happens once $k < 1/3$
(Section~\ref{sec:kirsch-appendix}). Rock has little tensile strength
and jointed rock effectively has none, so a
computed tensile zone means that the crown is destressed, not that it
carries tension. Joints open, blocks lose clamping, and structurally
controlled failure becomes possible in ground that the stress analysis
describes as understressed.

The failure zone is delineated by comparing the computed stress field
against a strength criterion point by point and mapping the region where
strength is exceeded. Its extent sets the
minimum length of reinforcement, since a bolt must be anchored beyond the
yielded rock. In brittle rock the depth of
spalling is better predicted by the empirical relations of
Chapter~\ref{ch:high-stress}, which capture behaviour that conventional criteria
do not.

\subsection{Excavation geometry and orientation}
\label{sec:geometry-orientation}

Shape controls the boundary stress distribution. Elongating an opening in the
direction of the major principal stress reduces the peak concentration; a
shape whose axis ratio matches the stress ratio approaches uniform boundary
stress and is the theoretical optimum (Section~\ref{sec:elliptical}). Sharp corners concentrate stress
severely, which is why square-profile openings are radiused, and flat roofs
are avoided in high horizontal stress for the same reason.

Size acts differently on the stress and on its consequences. In an infinite homogeneous elastic field
the concentration factor does not depend on span, so increasing the span does
not raise the tangential stress. This independence fails once the span
becomes comparable with the cover. What the span does increase is the volume of overstressed rock, the
weight of any block that forms, and the probability that an adverse structure
is intersected at all. The last of these is often the dominant effect, and it
is why experience from small openings extrapolates poorly to large ones.

Orientation has to satisfy two requirements. Relative to stress, the favourable
alignment places the long axis parallel to the major horizontal stress.
Relative to structure, the favourable alignment avoids running parallel to a
persistent set and minimises adverse wedge geometry
(Chapter~\ref{ch:structural-analysis}). The two rarely coincide, and the
alignment is then set by the mechanism that governs: stress where high stress
is combined with sparse structure, and structure where moderate stress is
combined with strong persistent jointing.

\section{Stratified and blocky rock}
\label{sec:stratified-and-blocky}

\subsection{Stratified rock and roof beams}
\label{sec:stratified-rock}

A bedded roof does not behave as a continuum. Bedding planes have little
tensile strength, so the immediate roof separates into layers that deflect
independently, and the design problem becomes one of beam or plate stability
rather than stress concentration.

The controlling variables are bed thickness and span. A thick bed spans
further than a thin one, and where several thin beds lie above the opening
they deflect together only if bonded or bolted; unbonded, the lowest bed
carries its own weight alone and fails first. Bed separation observed by
extensometer is therefore an early and reliable indicator of roof
deterioration.

A jointed or fractured roof bed cannot carry tension and does not act as an
elastic beam. It acts instead as a linear arch (the voussoir beam~\citep{beer1983voussoir}),
developing a compressive arch between abutments, with the thrust line rising
towards midspan, where contact is in the upper part of the bed, and falling
towards the abutments, where it is in the lower part
(Figure~\ref{fig:voussoir-roof}). Three failure modes
follow: crushing where the compressive stress at those contact zones exceeds
the rock strength, buckling or snap-through where the arch becomes too
shallow to be stable, and shear or slip at the abutments. The solution is iterative, since the moment arm depends on the
thrust and the thrust on the moment arm, and it yields a maximum stable span
that falls sharply as bed thickness decreases.

\begin{figure}[htbp]
\centering
\datafigure[\textwidth]{ch13/voussoir-roof}
\caption{Voussoir beam in a jointed roof bed of span $s$ and thickness $t$ (schematic, deflection exaggerated). The thrust line rises from the lower part of the bed at the abutments to the upper part at midspan, and compression concentrates in the contact zones, where crushing is checked.}
\label{fig:voussoir-roof}
\end{figure}

Large-deflection behaviour distinguishes a roof that sags from one that fails.
A voussoir beam can deflect substantially while remaining stable, and
continued deflection under time-dependent loading may either converge to a
stable geometry or accelerate into snap-through. The two are separated by
monitoring the rate of deflection rather than its magnitude.

\subsection{Support of laminated roofs}
\label{sec:laminated-roof-support}

Bolts in a bedded roof act in one of three different ways, and the
design has to choose among them.

\emph{Suspension} hangs the immediate weak beds from a competent horizon
above. It requires that such a horizon exists within reach and that the bolt
can be anchored in it, and the bolt load is then the weight suspended:
a 1.5~m by 1.5~m pattern carrying 1.0~m of laminated rock of unit weight
25~kN/m\textsuperscript{3} loads each bolt with 56~kN
(Figure~\ref{fig:roof-bolt-load}).
\emph{Beam building} clamps several thin beds together so they act as one
thick beam. Beam stiffness rises with the cube of
thickness, so four beds bolted into one beam are sixteen times as stiff
as the same four acting separately. Beam building requires closely spaced bolts and pretension sufficient to develop friction across the
bedding. \emph{Clamping} increases normal stress
across the bedding to prevent slip at the abutments, one of the three ways a
voussoir arch fails.

\begin{figure}[htbp]
\centering
\datafigure[\textwidth]{ch13/roof-bolt-load}
\caption{Left: suspension load per bolt, $W = \gamma h s_b^2$, for suspended thickness $h$ on a square pattern of spacing $s_b$ ($\gamma = 25$~kN/m\textsuperscript{3}); the point is the 56~kN of the text and the band the yield load of a 20~mm rebar (Table~\ref{tab:support-props}). Right: stiffness of $n$ beds bolted into one beam relative to the same beds acting separately, $n^2$.}
\label{fig:roof-bolt-load}
\end{figure}

Selection follows from bed thickness and available anchorage. Where a
competent horizon lies within bolt length, suspension is the most reliable and
the easiest to verify. Where it does not, bolts alone must build a beam,
and its effectiveness depends on pretension being achieved and
maintained, which is a matter of installation quality rather than of design.
Clamping governs where the mode of concern is slip at the abutments of a
voussoir bed, and where no competent horizon lies within the length of a
bolt, cable support extends the reach (Chapter~\ref{ch:openings}).

\subsection{Blocky rock}
\label{sec:blocky-rock}

In blocky ground the intact rock is strong and the discontinuities are weak
and persistent, so the rock mass fails by blocks moving rather than by
material yielding. The identification of those blocks is geometric and is
treated in Chapter~\ref{ch:structural-analysis}; what concerns this chapter is
their stability once identified.

The analysis is by limit equilibrium. Weight drives the block, friction, cohesion and
clamping stress on the sliding planes resist it, and support makes up the deficit.
Three points recur. First, clamping stress is the largest resisting contribution in a
confined crown and is nearly absent in a relaxed sidewall, so the same block
geometry is far more critical on a wall than in a roof. Second, water pressure on the
sliding planes reduces the effective normal stress and can halve the available resistance.
Third, the factor of safety is highly sensitive to the friction angle assumed for
a surface that has not been tested, which is the largest uncertainty in most such
calculations.

Key block theory~\citep{goodman1985block} offers a design economy. If the critical key block, the
one whose removal permits others to follow, is stabilised, the blocks it
keys are secured with it. A systematic bolting pattern designed against the
critical key block therefore covers a family of lesser blocks without each
being analysed, which is the technical justification for pattern bolting.

\section{Fractured rock, soft ground and the analysis route}
\label{sec:weak-ground}

\subsection{Sidewalls and heavily fractured rock}
\label{sec:sidewalls-fractured}

Tall sidewalls in high horizontal stress present a distinct problem. Stress
relief on excavation reduces confinement normal to the wall, steeply dipping
foliation or bedding can buckle into the opening, and blocks that would be
clamped in a crown are free to slide. Sidewall support therefore has to be designed for its own mechanism
and cannot be scaled down from the crown support. In large caverns
the walls frequently require more reinforcement than the roof.

As discontinuity spacing falls relative to the opening, block-based analysis
loses meaning. When an opening intersects hundreds of blocks rather than a
handful, no individual block governs and the mass again behaves as a
continuum, now weak and highly deformable. The transition is gradual and the
practical criterion is the ratio of block size to opening span, the $l_b/B$
of Section~\ref{sec:behaviour-classes}: where
that ratio falls below roughly a tenth, equivalent continuum treatment with
Hoek--Brown parameters (Chapter~\ref{ch:rock-mass-parameters}) is appropriate,
and where it exceeds about a third, discrete analysis is required. The two
values are indicative guides to the choice of analysis and not limits. Between them both analyses should be run;
disagreement usually means that the two represent different mechanisms
(Section~\ref{sec:methods-together}).

\subsection{Soft ground behaviour}
\label{sec:soft-ground-behaviour}

Soft ground fails through the material. The behaviour descriptions of
Terzaghi's tunnelman's ground classification~\citep{terzaghi1950geologic}, in
the form restated by \citet{heuer1974important} (firm, ravelling, running,
flowing, squeezing and swelling),
remain the most useful vocabulary because each implies a different control
measure. Firm ground stands unsupported long enough to install support;
ravelling ground loosens progressively from the face; running ground flows
until it reaches its angle of repose; flowing ground moves as a liquid and
requires face pressure; squeezing ground closes on the excavation; swelling
ground expands on contact with water.

Whether the response is drained or undrained governs the analysis. In
low-permeability clay, excavation is effectively undrained in the short term
and the total stress analysis applies to face stability; drainage over
subsequent years transfers load to the lining, so the long-term case is a
different calculation. In permeable ground the response is drained
throughout. The intermediate case (partial drainage at a rate comparable to
the advance rate) is common in weathered rock and silty soils and is the
hardest to analyse; conservatively, both bounds are checked.

Stand-up time links behaviour to construction. It is the period for which an unsupported
span remains stable; it falls sharply with span, and it determines
whether cyclic excavation is feasible. Estimates follow from the RMR
chart of span against stand-up time (Chapter~\ref{ch:empirical-design}),
whose values are indicative only. Where stand-up time approaches
the cycle time, the choice is to reduce the span through sequential
excavation, to improve the ground, or to change to a method that supports the
face continuously.

\subsection{Selecting the analysis route}
\label{sec:analysis-route}

The behaviour class determines the analysis. The mapping is stated
explicitly here because the commonest error in
underground analysis is a correct calculation of
the wrong mechanism.

Massive rock at moderate stress is analysed elastically, with
boundary stress checks (Chapter~\ref{ch:ccm}). Massive rock at
high stress requires spalling prediction and depth-of-failure estimation
(Chapter~\ref{ch:high-stress}). Stratified rock is analysed as a beam or voussoir arch
(Section~\ref{sec:stratified-and-blocky}), and blocky rock by structural analysis and
limit equilibrium (Chapter~\ref{ch:structural-analysis}). Heavily fractured rock is treated as an
equivalent continuum with empirical or numerical methods
(Chapters~\ref{ch:empirical-design} and~\ref{ch:numerical-modelling}). Weak
ground is analysed for face stability and deformation control
(Sections~\ref{sec:soft-ground-behaviour} and~\ref{sec:face-support}). The mapping is not a substitute
for judgement, but a design that applies the wrong route has usually done so
by habit rather than by decision.

\section{Face support and ground loss}
\label{sec:face-support}

\subsection{Face stability analysis}
\label{sec:face-stability-analysis}

The oldest treatment expresses face stability through the stability ratio,
which compares the total overburden and any surcharge, less the support
pressure applied at the face, against the undrained shear strength of the
ground,
\[
N = \frac{\sigma_s + \gamma z - \sigma_T}{s_u} ,
\]
where $\sigma_s$ is the surface surcharge, $\gamma z$ the total overburden
stress at the level of the tunnel axis ($\gamma$ the bulk unit weight of the
ground in kN/m\textsuperscript{3} and $z$ the depth of the axis in m),
$\sigma_T$ the support pressure applied at the face and $s_u$ the undrained
shear strength, the stresses all in kPa, so that $N$ is dimensionless. The observation of
\citet{broms1967stability} that instability occurs above a ratio of about six
remains a useful screening test. For an unsurcharged, unsupported face with
$\gamma = 20$~kN/m\textsuperscript{3} at an axis depth of 15~m, the overburden
stress is $20 \times 15 = 300$~kPa, and $N \le 6$ requires
$s_u \ge 300/6 = 50$~kPa; softer ground needs face pressure or
improvement (Figure~\ref{fig:stability-ratio-pressure}). \citet{davis1980stability}
refined the ratio into upper and lower bound plasticity solutions in
which the critical ratio depends on the cover-to-diameter ratio and on the
length of unsupported heading.

\begin{figure}[htbp]
\centering
\datafigure[0.62\textwidth]{ch13/stability-ratio-pressure}
\caption{Face pressure for a stability ratio of six, $\sigma_T = \gamma z - 6 s_u$, with $\gamma = 20$~kN/m\textsuperscript{3} and no surcharge. The point is the example of the text: at $z = 15$~m the face needs no pressure once $s_u \ge 50$~kPa.}
\label{fig:stability-ratio-pressure}
\end{figure}

These solutions are for cohesive ground and total stress. Frictional ground
requires a different construction, and the standard one is the wedge--silo
model~\citep{anagnostou1996face}: a wedge of ground slides forward at the face, loaded from above by a
prism whose weight is reduced by arching within it, calculated by Janssen's
silo theory~\citep{janssen1895versuche}. The required support pressure follows from equilibrium of the
wedge.

Two results from this model matter for design
(Figure~\ref{fig:wedge-silo}). The first is that the required support pressure
is insensitive to cover once the cover is sufficient for arching to develop,
so a deeper tunnel does not need proportionally more face pressure,
contrary to what reasoning from overburden suggests. At shallower cover
the weight of the prism acts nearly in full, and the required pressure does
rise with cover. The second is that the required pressure is highly sensitive to
cohesion: a small effective cohesion greatly reduces it,
which is why an apparently modest ground improvement can greatly improve face
stability, and why assuming zero cohesion in a slightly cemented soil is
severely conservative. The same sensitivity makes cohesion the parameter
least safely assumed: a small cemented component is easily destroyed by
disturbance and is rarely continuous along an alignment, so credit taken for
it is supported by testing and by a check at zero cohesion.

\begin{figure}[htbp]
\centering
\datafigure[\textwidth]{ch13/wedge-silo}
\caption{Wedge--silo model~\citep{anagnostou1996face}. Left: the wedge at the face, loaded by a prism whose weight is reduced by shear on its sides (schematic). Right: required face pressure for a square face of side 10~m in dry ground ($\gamma = 20$~kN/m\textsuperscript{3}, $\phi' = 30^\circ$); it is constant once $C/D$ exceeds about one and falls sharply with cohesion.}
\label{fig:wedge-silo}
\end{figure}

\subsection{Face support pressure}
\label{sec:face-support-pressure}

Closed-face machines apply pressure to the face, and the design task is to
define the operating window. The lower bound is the pressure preventing
collapse, from the wedge--silo analysis above, whose equilibrium is
formulated in full by \citet{anagnostou1996face}, plus the water pressure at
the face; the
upper bound is the pressure causing blow-out or heave, governed by the weight
of overburden and reducing sharply as cover decreases. In shallow ground the
window can be narrow (Figure~\ref{fig:face-pressure-window}), and in very shallow ground under a river or a soft
deposit it can close altogether, which is a constraint on the alignment that
the choice or operation of the machine cannot remove.

\begin{figure}[htbp]
\centering
\datafigure[0.62\textwidth]{ch13/face-pressure-window}
\caption{Operating window at the crown of a 10~m face below a water table at the surface. Collapse: wedge--silo pressure in submerged ground ($\phi' = 30^\circ$, $c' = 0$), taken as uniform over the face, plus pore pressure. Blow-out: total overburden, $\gamma = 20$~kN/m\textsuperscript{3}. The window closes at a cover of about 1.6~m.}
\label{fig:face-pressure-window}
\end{figure}

Pressure varies over the face height with the ground and slurry densities, so
the specification must state the single reference level at which the target
applies, conventionally the crown or the axis. The pressure gradient of the
chamber medium is lower than that of the overburden, so the margin against
blow-out is smallest at the crown. Where the margin against collapse is
smallest depends on how the required pressure is taken to vary. If the
effective part is uniform over the face, as in the wedge--silo resultant of
Figure~\ref{fig:face-pressure-window}, a slurry denser than water gains on the
pore pressure with depth and the collapse margin is also smallest at the crown.
If the effective part is taken to rise with depth as an earth pressure, the
demand can rise faster than the chamber pressure and the collapse margin is
smallest at the invert (Chapter~\ref{ch:tbm}). Each bound is therefore
evaluated at the elevation where it governs.

How pressure is transferred depends on the machine. A slurry shield transfers
pressure through a filter cake at the face, and its formation requires the
slurry to be able to bridge the pore openings; in coarse ground the slurry
penetrates instead of forming a cake, and the support pressure is lost into
the ground. An earth pressure balance machine transfers pressure through the
excavated soil itself, held as a plastic paste in the chamber, which requires
conditioning with foam or polymer to give material of the right consistency;
outside its grading envelope, conditioning becomes expensive or ineffective.
Machine selection is therefore governed by the ground before it is governed by plant
(Chapter~\ref{ch:tbm}).

\subsection{Ground loss and mixed face conditions}
\label{sec:ground-loss}

Volume loss is the sum of three separable
contributions~\citep{mair1996settlement}, which are the
construction-controlled part of the sources decomposed in
Chapter~\ref{ch:numerical-modelling}. \emph{Face loss} arises from movement into the face when
support pressure is below the in-situ stress. \emph{Shield loss} arises over
the length of the shield, from the overcut, from ploughing as the machine
steers, and from any pitch or yaw. \emph{Tail loss} arises at the shield tail,
where the annular gap between the excavated bore and the segment ring must be
filled by grout before the ground converges into it.

Separating them is useful because each has a different control.
Face loss is controlled by chamber pressure, shield loss by overcut and
steering discipline, and tail loss by grout pressure, volume and timing. A
measured volume loss above target cannot be acted on until it is attributed
to one of the three; the attribution identifies the control to be changed. Tail loss is usually the
largest contributor and the most controllable.

Mixed-face conditions (rock in the invert and soil in the crown, or the
reverse) are among the most damaging situations in mechanised tunnelling.
The tools wear unevenly, the machine tends to steer towards the softer
material, the softer horizon is over-excavated, and face pressure appropriate
to one material is wrong for the other. Volume loss in mixed-face conditions
commonly exceeds the design assumption. Since these conditions arise wherever
a strength contrast crosses the face, most often an irregular rockhead
(Chapter~\ref{ch:engineering-geology}) but also faults, weathering fronts,
dykes and boulder horizons, predicting where they occur is a matter of
ground modelling and of investigation designed to resolve that contrast, not
of machine operation.

\section{Alternatives to face pressure}
\label{sec:alternatives-face-pressure}

Supporting the face is one response to weak ground; changing the ground is
another, and sometimes the better one. Where the required face pressure is
unachievable, where the operating window has closed, or where settlement
limits are tighter than pressure control can deliver, treating the ground
ahead of the face converts the problem into one the excavation method can
handle.

The options are set out in Chapter~\ref{ch:ground-improvement} and span
grouting, jet grouting, freezing, dewatering and pre-support by umbrella
arches or face dowels. Their common feature is a cost in programme as well
as in money: treatment must be completed before excavation reaches it, so
it either occupies the critical path or must be executed from the surface in
advance.

The judgement is therefore economic as much as technical. Continuous face
support carries a recurring cost in conditioning agents or slurry treatment
and in advance rate, and constrains what ground the drive can pass through,
but it is available along the whole drive; ground treatment costs time and
money at specific locations and removes the constraint there. Long drives in mostly
competent ground with a few difficult reaches favour treatment at those
reaches; drives predominantly in weak ground favour a machine that supports
the face throughout. The decision is taken at the same time as machine
selection, and reversing it later is rarely possible.

\chapterreferences
